\chapter*{Abstract}
\addcontentsline{toc}{chapter}{Abstract}

Formal specifications make software behaviour machine-checkable, yet they are
rarely written in practice because authoring them by hand is costly and demands
expertise most development teams lack. The machinery for \emph{checking} programs
against contracts is mature; what has been missing is an economical supply of the
contracts. This thesis argues that the obstacle is not the value of specifications
but their cost of acquisition, and that specifications inferred automatically from
code---once cheap enough to be useful---can be composed to reason about library
compatibility and integrated as the verification layer of the AI-native
development lifecycles now emerging around large language models.

The work is built on a single instrument: the JML-Inferrer, a static analysis
tool that derives Java Modeling Language contracts---preconditions,
postconditions, loop and class invariants, and frame conditions---from Java source
by heuristics over the abstract syntax tree, and holds them to the demanding
standard of discharge by the OpenJML extended static checker and the Z3 SMT
solver. The thesis develops three connected results. First, it establishes that
heuristic, partial, automatically inferred specifications are of high enough
fidelity to be useful, and that supplying them to a large language model
measurably improves the unit tests the model generates, converting the abstract
value of a contract into a measured downstream benefit. Second, it shows that a
compositional weakest-precondition pass over the call graph strictly extends
single-pass inference, recovering a substantial number of additional well-formed
preconditions and enabling behavioural version-compatibility to be treated as a
static, run-free analysis: a change to a callee's contract can be propagated to
its callers to compute the blast radius of the change, turning semantic
versioning from an asserted convention into a checkable property. Third, it shows
how inferred, machine-checkable contracts supply what current agentic,
specification-driven lifecycles lack---a contract layer that is machine-checkable
rather than natural language, and a closed verification loop rather than an open
generate-and-review cycle---with the same contract serving as both the input an
agent builds against and the oracle its output is checked against.

Throughout, the thesis is scrupulous to distinguish \emph{verified}
specifications, which the checker has discharged, from those that are merely
\emph{well-formed}. Taken together, the results trace a single arc from
inference---making specifications cheap---to integration---making them
load-bearing in how software is built.
